\begin{table*}[t]
\centering
\small
\renewcommand{\arraystretch}{1.25} % Adds breathing room between rows
\setlength{\tabcolsep}{10pt} % Adds space between columns

\resizebox{0.95\textwidth}{!}{%
\begin{tabular}{llll}
\toprule
\textbf{Abbr.} & \textbf{Task} & \textbf{Dataset} & \textbf{Lang.} \\
\midrule

\rowcolor{colorint!10}
\multicolumn{4}{l}{\textbf{Intrinsic: Core Capability} (Metrics $\downarrow$ Lower is better)} \\
\addlinespace[0.1cm]

\multicolumn{4}{l}{\hspace{0.3cm}\textit{Phone Recognition (PFER)}} \\
\texttt{PR-tmt} & Variation of Seen Language & TIMIT \cite{garofolo1993timit} & English \\
\texttt{PR-arc} & Variation of Seen Language & L2-ARCTIC Perceived \cite{zhao2018l2arctic} & English \\
\texttt{PR-saa} & Variation of Seen Language & Speech Accent Archive \cite{gmuspeechaccentarchive} & English \\
\texttt{PR-drc} & Unseen Languages & DoReCo \cite{paschen2020doreco} & 45 langs \\
\texttt{PR-vox} & Unseen Languages & VoxAngeles \cite{chodroff2024voxangeles} & 95 langs \\
\texttt{PR-tsm} & Unseen Languages & Tusom2021 \cite{mortensen2021tusom2021} & Tusom \\

\addlinespace[0.1cm]
\midrule
\rowcolor{colorext!10}
\multicolumn{4}{l}{\textbf{Extrinsic: Downstream Utility} (Metrics $\uparrow$ Higher is better)} \\
\addlinespace[0.1cm]

\multicolumn{4}{l}{\hspace{0.3cm}\textit{Pathological Speech: Dysarthria Intelligibility Prediction ($\tau$) \& Child Speech Disorder Detection (F1)}} \\
\texttt{DYS-ez} & Dysarthria Intelligibility Prediction & EasyCall \cite{turrisi21_interspeech} & Italian \\
\texttt{DYS-ua} & Dysarthria Intelligibility Prediction & UASpeech \cite{kim08c_interspeech} & English \\
\texttt{CSD-us} & Child Speech Disorder Detection & UltraSuite \cite{eshky2018ultrasuite} & English \\

\addlinespace[0.2cm]
\multicolumn{4}{l}{\hspace{0.3cm}\textit{L2 Speech: L1 Classification (F1) \& L2 Assessment ($\tau$)}} \\
\texttt{L1-eda} & L1 Classification & EdAcc \cite{sanabria2023edacc} & English \\
\texttt{L1-arc} & L1 Classification & \citet{kominek04_ssw} \& \citet{zhao2018l2arctic} & English\\
\texttt{L2-so} & L2 Assessment & Speechocean762 \cite{zhang2021speechocean762} & English\\

\addlinespace[0.2cm]
\multicolumn{4}{l}{\hspace{0.3cm}\textit{Multilingual: Lang. ID (F1), Geolocation (Recall@1) \& Phone Inventory Induction (F1-PI)}} \\
\texttt{LID-fl} & Lang. ID (LID) & FLEURS-24 \cite{conneau2023fleurs} & 24 langs \\
\texttt{GEO-v}  & Speech Geolocation & Vaani \cite{vaani2025} & Hindi Dialects \\
\texttt{PI-drc} & Phone Inventory Induction & DoReCo \cite{paschen2020doreco} & 45 langs \\

\bottomrule
\end{tabular}
}
\caption{
List of evaluation tasks. \colorbox{colorint!10}{Blue} denotes core capabilities, where lower scores are better. \colorbox{colorext!10}{Yellow} denotes downstream utility, where higher scores are better. \textit{F1-PI} is described in \autoref{appendix:phoneinv}. See \autoref{appendix:datadetails} for license details.
}
\label{tab:evaltasks}
\end{table*}


\section{Background}

\subsection{Phone Recognition Systems}\label{ss:bg-pr-system}
PR can be viewed as a variant of the S2T task that maps speech to phonetic symbols such as IPA \cite{ipa1999handbook}. In this work, we use ``PR system'' to refer broadly to any system capable of converting speech into IPA in a language-agnostic fashion.

Modern PR systems are typically fine-tuned from ASR systems \cite{baevski2020wav2vec, radford2023robust} or trained from scratch on ASR datasets \cite{zhu-etal-2024-taste} with transcriptions automatically converted to IPA using grapheme-to-phoneme (G2P) tools \cite{mortensen2018epitran, zhu2022charsiu-g2p}.
Language-specific approaches \cite{li2020universal, gao21_interspeech} rely on phoneme inventories, while language-agnostic approaches, which we focus on, seek to learn phonetic representations generalized across languages.
LALMs have recently become prominent in speech tasks and have shown competitive performance with cascaded systems that combine LLMs with speech processing modules \cite{yang2025holisticlalm}, motivating interest in their application to PR \cite{huang2025dynamic, wang2025audiobench}.
We describe the systems investigated in this work in \autoref{ss:exp_models}.


\subsection{Phonetic information in PR systems}\label{ss:bg-phonetic-info}
Explicitly generated phonetic transcriptions are easy for humans to inspect and utilize. 
For example, faithful phonetic transcriptions of the speech of a child with a speech sound disorder can help a clinician understand the nature of the disorder and design interventions \cite{dodd2013differential}.
Nevertheless, representing continuous speech with discrete symbols inherently incurs information loss, filtering out non-linguistic variation.

Internal model representations serve as a complement that retains richer information.
Speech models trained on S2T tasks produce temporally aligned representations that capture empirically useful acoustic-phonetic \cite{choi24b_s3mphonetic}, articulatory \cite{cho2024sslartic}, and even semantic \cite{ma-etal-2025-cross} features.
The most widely used S2T model representations are from end-to-end ASR models such as Whisper \cite{radford2023robust} and WavLM \cite{chen2022wavlm}. 
In contrast, LALMs’ representations are often inaccessible or difficult to analyze, as they focus mainly on textual output and lack strict temporal alignment with input speech.


\subsection{Assessing phonetic/phonological ability}
In the text modality, language models are evaluated with text input and output. 
PhonologyBench \cite{suvarna2024phonologybench} evaluates G2P, syllable counting, and rhyme judgment, while \citet{bunzeck2025smalllm} and \citet{goriely2025wordseg} probe phonological knowledge using minimal pairs and word segmentation.
In the speech modality, models are evaluated with speech (and optionally text) input and representation output. 
SUPERB \cite{yang21c_interspeech} and Dynamic-SUPERB \cite{huang2025dynamic} include phoneme recognition, phonological feature analysis, and pronunciation evaluation.
BabySLM \cite{lavechin2023babyslm} and the ZeroSpeech challenges \cite{nguyen2020zerospeech} propose metrics that evaluate phonological and acoustic-phonetic contrasts based on minimal pairs.
In contrast to previous work, \bench evaluates phonetic ability in both text and speech through intrinsic and extrinsic tasks. 